\section*{Розділ VII. Комісії та робочі органи КСУ}
\addcontentsline{toc}{section}{Розділ VII. Комісії та робочі органи КСУ}

\subsection*{7.1. Створення тимчасових комісій}
\addcontentsline{toc}{subsection}{7.1. Створення тимчасових комісій}
    7.1.1. КСУ має право створювати тимчасові комісії для виконання окремих завдань, підготовки питань до розгляду або здійснення спеціалізованих функцій.

    7.1.2. Ініціатива про створення тимчасової комісії може бути внесена:

        \begin{enumerate}[label=\alph*)]
            \item Групою делегатів КСУ кількістю не менше 3 осіб;
            \item Спікером КСУ;
            \item Органом студентського самоврядування через своїх делегатів у КСУ;
            \item За дорученням КСУ у ході розгляду конкретного питання.
        \end{enumerate}

    7.1.3. Рішення про створення тимчасової комісії приймається КСУ простою більшістю голосів від присутніх делегатів та повинно містити:

        \begin{enumerate}[label=\alph*)]
            \item Назву комісії;
            \item Мету та завдання комісії;
            \item Орієнтовний склад комісії (кількість членів);
            \item Строки роботи комісії;
            \item Повноваження комісії;
            \item Порядок звітування комісії перед КСУ.
        \end{enumerate}

\subsection*{7.2. Порядок формування складу комісій}
\addcontentsline{toc}{subsection}{7.2. Порядок формування складу комісій}
    7.2.1. До складу тимчасових комісій КСУ можуть входити:

        \begin{enumerate}[label=\alph*)]
            \item Делегати КСУ;
            \item Представники органів студентського самоврядування (за згодою відповідного органу);
            \item Експерти та консультанти з питань, що належать до компетенції комісії;
            \item Представники студентських організацій (за їх згодою).
        \end{enumerate}

    7.2.2. Склад комісії формується за принципом добровільності участі та професійної компетентності у відповідній сфері.

    7.2.3. Голова комісії обирається з числа її членів простою більшістю голосів членів комісії на її першому засіданні або призначається рішенням КСУ про створення комісії.

    7.2.4. Кількісний склад комісії визначається характером та обсягом покладених на неї завдань, але не може бути менше 3 осіб.

    7.2.5. Члени комісії мають рівні права у її роботі, незалежно від їх статусу в системі студентського самоврядування.

\subsection*{7.3. Повноваження комісій}
\addcontentsline{toc}{subsection}{7.3. Повноваження комісій}
    7.3.1. Тимчасові комісії КСУ у межах покладених на них завдань мають право:

        \begin{enumerate}[label=\alph*)]
            \item Збирати та аналізувати інформацію з питань своєї компетенції;
            \item Підготовлювати проєкти рішень КСУ;
            \item Запитувати необхідні документи та матеріали від органів студентського самоврядування;
            \item Проводити власні засідання та робочі зустрічі;
            \item Заслуховувати звіти та інформацію від посадових осіб органів студентського самоврядування;
            \item Запрошувати на свої засідання представників адміністрації Університету, експертів та інших осіб;
            \item Створювати робочі групи з числа своїх членів для опрацювання окремих питань.
        \end{enumerate}

    7.3.2. Органи студентського самоврядування зобов'язані сприяти роботі комісій КСУ, надавати запитувану інформацію та документи у строки, встановлені комісією.

    7.3.3. Комісії КСУ не мають права приймати рішення, обов'язкові для виконання органами студентського самоврядування, окрім випадків, прямо передбачених рішенням КСУ про створення комісії.

    7.3.4. Рекомендації та висновки комісій мають консультативний характер та підлягають розгляду КСУ для прийняття остаточного рішення.

\subsection*{7.4. Порядок роботи комісій}
\addcontentsline{toc}{subsection}{7.4. Порядок роботи комісій}
    7.4.1. Комісії КСУ працюють на засадах колегіальності, відкритості та ефективності.

    7.4.2. Засідання комісії скликаються її головою за власною ініціативою або за вимогою не менше третини членів комісії.

    7.4.3. Засідання комісії є правомочним, якщо на ньому присутні не менше половини від загального складу членів комісії.

    7.4.4. Рішення комісії приймаються простою більшістю голосів від присутніх членів. У разі рівності голосів, голос голови комісії є вирішальним.

    7.4.5. Комісія може проводити засідання у дистанційному форматі з використанням засобів електронного зв'язку за аналогією з процедурами, встановленими для КСУ.

    7.4.6. Хід засідань комісії фіксується у протоколі, який підписується головою комісії та секретарем засідання.

    7.4.7. Комісія визначає свій внутрішній регламент роботи, який не може суперечити цьому Регламенту та рішенню КСУ про створення комісії.

\subsection*{7.5. Строки роботи комісій}
\addcontentsline{toc}{subsection}{7.5. Строки роботи комісій}
    7.5.1. Строки роботи тимчасової комісії визначаються рішенням КСУ про її створення залежно від характеру та складності покладених завдань.

    7.5.2. Строки роботи комісії можуть визначатися:

        \begin{enumerate}[label=\alph*)]
            \item Конкретною датою завершення роботи;
            \item Моментом виконання покладених завдань;
            \item Прив'язкою до певної події (наприклад, до наступного засідання КСУ);
            \item Строком повноважень поточного складу КСУ.
        \end{enumerate}

    7.5.3. КСУ може продовжити строки роботи комісії за її мотивованим поданням або достроково припинити її діяльність у разі виконання завдань або зміни обставин.

    7.5.4. Комісія зобов'язана завершити свою роботу не пізніше завершення повноважень складу КСУ, який її створив. У разі обрання нового складу КСУ діяльність комісії припиняється, якщо новий склад КСУ на першому засіданні не прийме рішення про продовження її роботи.

\subsection*{7.6. Звітність комісій}
\addcontentsline{toc}{subsection}{7.6. Звітність комісій}
    7.6.1. Комісії КСУ є підзвітними КСУ та зобов'язані регулярно інформувати про хід виконання покладених завдань.

    7.6.2. Періодичність звітування визначається рішенням КСУ про створення комісії, але не може бути рідше одного разу на три місяці.

    7.6.3. Підсумковий звіт комісії подається КСУ після завершення основних завдань та містить:

        \begin{enumerate}[label=\alph*)]
            \item Аналіз виконаної роботи;
            \item Основні висновки та рекомендації;
            \item Проєкти рішень КСУ (за потреби);
            \item Пропозиції щодо подальших дій;
            \item Інформацію про використані ресурси та витрачений час.
        \end{enumerate}

    7.6.4. КСУ розглядає звіти комісій на своїх засіданнях та приймає рішення щодо запропонованих рекомендацій.

\subsection*{7.7. Спеціалізовані комісії}
\addcontentsline{toc}{subsection}{7.7. Спеціалізовані комісії}
    7.7.1. Для розгляду окремих категорій питань КСУ може створювати спеціалізовані комісії з особливими процедурами роботи:

        \begin{enumerate}[label=\alph*)]
            \item \textbf{Апеляційна комісія} -- для розгляду апеляцій на рішення СУД відповідно до Розділу V цього Регламенту;
            \item \textbf{Лічильна комісія} -- для підрахунку голосів при таємному голосуванні;
            \item \textbf{Комісія з питань регламенту} -- для підготовки пропозицій щодо вдосконалення процедур роботи КСУ.
        \end{enumerate}

    7.7.2. Особливості роботи спеціалізованих комісій визначаються відповідними розділами цього Регламенту або окремими рішеннями КСУ.

    7.7.3. До складу апеляційної комісії не можуть входити особи, які мають конфлікт інтересів у справі, що розглядається.

\subsection*{7.8. Припинення діяльності комісій}
\addcontentsline{toc}{subsection}{7.8. Припинення діяльності комісій}
    7.8.1. Діяльність тимчасової комісії КСУ припиняється:

        \begin{enumerate}[label=\alph*)]
            \item У зв'язку із закінченням строків, визначених рішенням про створення;
            \item Після виконання покладених завдань та подання підсумкового звіту;
            \item За рішенням КСУ про дострокове припинення діяльності;
            \item У зв'язку із завершенням повноважень складу КСУ.
        \end{enumerate}

    7.8.2. Комісія може ініціювати припинення власної діяльності за рішенням більшості її членів у разі неможливості виконання покладених завдань.

    7.8.3. Після припинення діяльності комісії всі її матеріали передаються до електронного архіву ОСС для зберігання.